\scanpage{277}
\mspage{7}

\deleted{\unsure{On}} \deleted{Pour voir qu'}\inserted{Par $\operatorname{Loc}(3'\,b)$} elle est dans $W$.

\deleted{\unsure{Considérons|Remarquons}}

\deleted{\[
\unsure{i|\iota}_X^{o}:X^{o}\longrightarrow\Phi_{S^{o}}(X^{o}),
\]}

\deleted{elle est $\in W$ par $\operatorname{Loc}(3'\,b)$, donc si on prouve}

\deleted{\[
W=W^{o},\tag*{$(\unsure{*})$}
\]}

\deleted{on en conclura $i_X\in W$. \deleted{\unsure{que}} Admettons pour le moment que $\operatorname{Loc}(1,2a,3')$ implique $(*)$, et terminons de prouver qu'\unsure{elles|ils} impliquent $\operatorname{Loc}(3)$. \deleted{$\unsure{\Psi_S}$} Appliquons $\operatorname{Loc}(3')$ à $(\Psi_S(X^{\deleted{\unsure{o}}}))^o\longrightarrow S^{o}$, i.e. à $\Phi_{S^{o}}(X^{o})\longrightarrow S^{o}$, dont les fibres sont les $X_s^{o}$. Comme les $X_s$ sont $W$-asphériques, il résulte de \unsure{$(*)$|ceci|car} qu'il en est de $\unsure{=}$ des $X_s^{o}$, donc $\operatorname{Loc}(3')$ implique que $\Phi_{S^{o}}(X^{o})\longrightarrow S^{\unsure{o}}$ est dans $W$, donc \unsure{utilisant} \deleted{à \unsure{nouveau} \unsure{$(*)$|car} on conclut que}}

d'autre part, les fibres \inserted{$ {}_{s\backslash}X$} de $\Phi_S(X)$ sont $W$-coasphériques par hypothèse, donc par $\operatorname{Loc}(3'\,a)$ $\Phi_S(X)\unsure{{-}|\longrightarrow}S$ est dans $W$. Comme $i_X\in W$, il s'ensuit $(f:X{-}S)\in W$, qed.

\[
\text{b bis)}\qquad\operatorname{Loc}(1,2a,3'\,\mathrm{bis})\Longrightarrow\operatorname{Loc}(1,2,3\,\mathrm{bis}).
\]

On \deleted{sait} a déjà prouvé $\operatorname{Loc}(1,2,3a)$ \deleted{$\unsure{(3)}$} (vu que $\operatorname{Loc}(3'\,\mathrm{bis})$ contient $\operatorname{Loc}(3')$ comme cas particulier), \deleted{prouvons} donc $\operatorname{Loc}(3\,\mathrm{bis})$ en admettant déjà $\unsure{\Longrightarrow|{-}}\,3a$. Il suffira même de le prouver sous forme duale, puisque

\[
\operatorname{Loc}(1,2,3\,\mathrm{bis})\Longleftrightarrow\operatorname{Loc}(1,2^{o},3\,\mathrm{bis}^{o}),
\]

% archive sheet number: 276
